\chapter{Лемма о факторизации}
\noindent В главе доказывается абстрактный критерий, который по суммируемым оценкам цены разрезов обеспечивает устойчивость корневых компонент вложенной фильтрации. Затем модуль путей и контактные графы связывают потерю локального радиуса с этой ценой. Для множества Мандельброта критерий даёт условную факторизацию при совместном выполнении гипотез AM и GC; сами оценки разбираются в следующих главах.

\section{Общие критерии через разрезы и модуль путей}
\label{sec:general-cut-theorems}

Сначала отделим общую схему от конкретной параметрической фильтрации,
вводимой ниже. Пусть
\((Z,d)\) — компактное метрическое пространство, а
\[
 Y_0\supseteq Y_1\supseteq\cdots
\]
— непустые компакты с \(X:=\bigcap_{n\ge0}Y_n\ne\varnothing\).
Обозначим через \(\overline B_d(x,t)\) замкнутый шар в \(Z\). Для
\(A\subseteq Z\), \(x\in A\) через \(\operatorname{Comp}_x(A)\)
обозначим компоненту связности \(A\), содержащую \(x\). Для
\(x\in X\), \(r>0\) положим
\[
 C_n^X(x,r):=
 \operatorname{Comp}_x\bigl(Y_n\cap\overline B_d(x,r)\bigr),\qquad
 \rho_n^X(x,r):=
 \sup\{\delta\in[0,r]:
 X\cap\overline B_d(x,\delta)\subseteq C_n^X(x,r)\}.
\]
Пусть \(\lambda_m\in\{0,1,2,\ldots\}\),
\(1\le\lambda_1<\lambda_2<\cdots\), и
\(\lambda_m\to\infty\). Сократим обозначения:
\[
 R_m:=\rho_{\lambda_m}^X(x,r),\qquad
 \Delta_m:=R_m-R_{m+1}\ge0.
\]
Для блока глубин разрешается задать любую неотрицательную величину
\(\mathfrak c_m\), интерпретируемую как цену конечного разделяющего
набора. В следующей теореме используются только её численные оценки.

\begin{theorem}[Абстрактный критерий суммируемой цены разрезов]
\label{thm:abstract-summable-cut}
Зафиксируем \(x\in X\), \(r>0\) и последовательность
\((\lambda_m)\), определённую выше. Пусть существуют числа
\(\mathfrak c_m\ge0\), \(\epsilon_m\ge0\) и константы \(A,B>0\),
не зависящие от \(m\), такие, что
\[
 \sum_m\epsilon_m<\infty,\qquad
 \exists m_0\ \forall m\ge m_0,\quad
 0\le\Delta_m\le A\mathfrak c_m,\qquad
 \mathfrak c_m\le B R_m\epsilon_m,\qquad R_m>0.
\]
Если множества \(Y_n\cap\overline B_d(x,r)\) имеют конечное число
компонент, то существует \(\delta>0\), для которого
\[
 X\cap\overline B_d(x,\delta)\subseteq C_n^X(x,r)
 \qquad(n\ge0).
\]
В частности, система компонент конечных приближений факторизуется
в смысле включения
\[
 \forall n\ \exists L\ge n,\qquad
 Y_L\cap\overline B_d(x,\delta)\subseteq C_n^X(x,r).
\]
Если для каждого \(x\in X\) и \(\varepsilon>0\) такие точечные данные
можно выбрать с \(r<\varepsilon\), то \(X\) локально связно.
\end{theorem}
\begin{proof}
Положим \(q_m:=AB\epsilon_m\). По предположению
\(\Delta_m\le q_mR_m\) при всех \(m\ge m_0\). Из суммируемости следует,
что \(q_m\le1/2\) для всех достаточно больших \(m\); увеличим \(m_0\),
чтобы это неравенство и \(R_m>0\) выполнялись при каждом \(m\ge m_0\).
Тогда
\[
 R_m\ge R_{m_0}\prod_{j=m_0}^{m-1}(1-q_j)
 \ge R_{m_0}\exp\!\left(-2\sum_{j=m_0}^{\infty}q_j\right)>0
 \qquad(m\ge m_0).
\]
Последовательность \(\rho_n^X(x,r)\) невозрастает, а
\(\lambda_m\to\infty\), поэтому
\(\inf_n\rho_n^X(x,r)=\inf_mR_m>0\). Выберем
\(0<\delta<\inf_n\rho_n^X(x,r)\). По определению супремума и
наследственности включения при уменьшении радиуса получаем
\[
 X\cap\overline B_d(x,\delta)\subseteq C_n^X(x,r)
 \qquad(n\ge0).
\]
Для фиксированного \(n\) множество
\[
 B_n:=\bigl(Y_n\cap\overline B_d(x,\delta)\bigr)
       \setminus C_n^X(x,r).
\]
Поскольку \(\delta<r\), это в точности пересечение
\(Y_n\cap\overline B_d(x,\delta)\) с объединением всех компонент
\(Y_n\cap\overline B_d(x,r)\), кроме \(C_n^X(x,r)\). Эти компоненты
компактны, и их конечность показывает, что \(B_n\) компактно.
Кроме того, \(B_n\cap X\) пусто. Если бы
\(Y_L\cap\overline B_d(x,\delta)\cap B_n\ne\varnothing\) для всех
\(L\ge n\), вложенность компактных множеств дала бы непустое
пересечение, лежащее в \(X\cap B_n\), противоречие. Следовательно,
для некоторого \(L\ge n\)
\[
 Y_L\cap\overline B_d(x,\delta)\subseteq C_n^X(x,r).
\]
Наконец,
\(C:=\bigcap_n C_n^X(x,r)\) — пересечение убывающих непустых
компактных связных множеств. Оно компактно и связно, содержится в
\(X\cap\overline B_d(x,r)\) и содержит
\(X\cap\overline B_d(x,\delta)\). Выбирая \(r<\varepsilon\), получаем
связную окрестность \(x\) внутри любого заданного относительного
окружения. Значит, \(X\) локально связно.
\end{proof}

\subsection{Две независимые оценки для геометрических фильтраций}

Связь цены разреза с потерей радиуса можно выразить без упоминания
динамики. Пусть \(\mathcal Q_m=\{Q_{m,\alpha}\}_{\alpha\in J_m}\) —
конечное семейство компактных ячеек, а \(\Gamma_m\) — семейство
непрерывных путей. Определим дискретный \(1\)-модуль
\[
 \operatorname{Mod}_{1,m}(\Gamma_m):=
 \inf\left\{
 \sum_{\alpha\in J_m}w_\alpha\operatorname{diam}Q_{m,\alpha}:
 \begin{array}{l}
 w_\alpha\ge0,\\[-1mm]
 \displaystyle
 \sum_{\alpha:\,Q_{m,\alpha}\cap\gamma([0,1])\ne\varnothing}
 w_\alpha\ge1\quad(\gamma\in\Gamma_m)
 \end{array}\right\}.
\]
 Далее рассматриваются блоки с \(\Delta_m>0\); при \(\Delta_m=0\)
 радиальное неравенство тривиально.
 Для блока с \(\Delta_m>0\) положим
\[
 \mathcal L_m^X(x,r):=
 X\cap\bigl(C_{\lambda_m}^X(x,r)
                \setminus C_{\lambda_{m+1}}^X(x,r)\bigr).
\]
Из определения радиуса следует
\[
 R_{m+1}=\min\{R_m,\dist(x,\mathcal L_m^X(x,r))\}.
\]
Действительно, при каждом \(t<R_m\) точки
\(X\cap\overline B_d(x,t)\) лежат в корневой компоненте уровня
\(\lambda_m\), а первые из них, исключаемые из корневой компоненты
уровня \(\lambda_{m+1}\), принадлежат
\(\mathcal L_m^X(x,r)\).
При конечном числе компонент среза
\(Y_{\lambda_{m+1}}\cap\overline B_d(x,r)\) множество
\(\mathcal L_m^X(x,r)\) является конечным объединением пересечений
этих компонент с \(X\cap C_{\lambda_m}^X(x,r)\); значит, оно компактно,
и существует ближайшая к \(x\) точка \(y_m\) в этом множестве.
Пусть
\[
 n_m:=\min\{n:\lambda_m<n\le\lambda_{m+1},\
                    y_m\notin C_n^X(x,r)\}.
\]
Предположим, что срез \(C_{n_m-1}^X(x,r)\) путём связен, и обозначим
через \(\Gamma_m\) семейство всех непрерывных путей в нём от \(x\) до
\(y_m\). Для компактных полуалгебраических срезов это свойство
следует из теоремы о триангуляции полуалгебраических множеств; ниже оно
применимо к конечным полиномиальным срезам. Пусть \(\mathcal Q_m\) —
конечное семейство ячеек положительного диаметра, покрывающее все интерфейсы блока,
включая \(\Sigma_{n_m}\), причём
\(\Sigma_{n_m}\) пересекает каждый путь из \(\Gamma_m\). Обозначим
через \(\widehat{\mathfrak c}_m\) минимальную цену графового разреза,
отделяющего корень от компоненты, содержащей \(y_m\); цена разреза,
представленного набором \(S\) ячеек, равна
\(\sum_{\alpha\in S}\operatorname{diam}Q_{m,\alpha}\).

Цена любого набора ячеек, пересекающего каждый путь из \(\Gamma_m\),
не меньше этого модуля: индикаторы выбранных ячеек допустимы в
инфимуме. В частности, если каждое реберное разделение в конечном
контактном графе пересекает все такие пути и
\(\widehat{\mathfrak c}_m\) — минимальная цена допустимого графового
разреза, то
\begin{equation}
 \operatorname{Mod}_{1,m}(\Gamma_m)\le\widehat{\mathfrak c}_m.
 \label{eq:modulus-below-cut}
\end{equation}
Следовательно, независимая от масштаба нижняя оценка
\begin{equation}
 \operatorname{Mod}_{1,m}(\Gamma_m)\ge a_{x,r}\,\Delta_m,\qquad
 a_{x,r}>0,
 \label{eq:annular-modulus-target}
\end{equation}
немедленно даёт радиальный переход
\(\Delta_m\le a_{x,r}^{-1}\widehat{\mathfrak c}_m\). Это формулировка в
языке дискретного модуля семейств кривых; её общий источник —
количественная толщина путей, а не только связность конечного графа.

Для планарной клеточной системы зададим контактную кратность
\[
 v_{m,\alpha}:=
 \#\{n:\lambda_m<n\le\lambda_{m+1},\
                 Q_{m,\alpha}\cap\Sigma_n\ne\varnothing\},
\]
а \(z_{m,\alpha}\in Q_{m,\alpha}\) выберем произвольно. Соответствующая
атомарная контактная мера и её полная масса равны
\[
 \Xi_m:=
 \sum_{\alpha\in J_m}
 v_{m,\alpha}\operatorname{diam}(Q_{m,\alpha})\delta_{z_{m,\alpha}}.
\]
Для определения цены разреза рассмотрим локальный срез
\(C_{n-1}^X(x,r)\), его внутреннюю часть
\(A_n:=C_{n-1}^X(x,r)\cap Y_n\) и внешнюю часть
\(U_n:=C_{n-1}^X(x,r)\setminus A_n\). Вершины конечного интерфейсного
графа — компоненты \(A_n\) и замыкания компонент \(U_n\); ребро означает
непустое пересечение. Для ребра \(e=(A,\overline U)\) контакт есть
\(P_e:=A\cap\overline U\). Предполагается, что \(\Sigma_n\) содержит
все эти контакты и что ячейки блока покрывают каждый интерфейс. Корневая
вершина — компонента \(A_n\), содержащая \(x\), то есть \(C_n^X(x,r)\).
Удаляться могут только рёбра между \(A_n\)-вершинами и
\(\overline U_n\)-вершинами; рёбра между двумя \(\overline U_n\)-вершинами
сохраняются. Обозначим через
\[
I_{n,m}:=\{\alpha\in J_m:Q_{m,\alpha}\cap\Sigma_n\ne\varnothing\}.
\]
Набор \(S\subseteq I_{n,m}\) удаляет ребро \(e=(A,\overline U)\), если
\(P_e\subseteq\bigcup_{\alpha\in S}Q_{m,\alpha}\); он допустим, если после
удаления корневая вершина не соединена ни с одной другой вершиной
\(A_n\), пересекающей \(X\). Пусть \(\operatorname{cap}_{n,m}\) —
минимальная цена допустимого разреза, где цена \(S\) равна
\(\sum_{\alpha\in S}\operatorname{diam}(Q_{m,\alpha})\). Положим
\[
\mathfrak c_m:=
\sum_{\lambda_m<n\le\lambda_{m+1}}\operatorname{cap}_{n,m}.
\]
Полный набор \(I_{n,m}\) покрывает все контакты, удаляет каждое ребро
между \(A_n\)- и \(\overline U_n\)-вершинами. Разные вершины,
заданные компонентами \(A_n\), не пересекаются, поэтому этот набор
допустим.
Перестановка конечных сумм даёт
\begin{equation}
 \mathfrak c_m\le\Xi_m(\C).
 \label{eq:block-capacity-contact-mass}
\end{equation}
Искомый контроль этой массы имеет естественную общую формулировку через
Carleson-упаковку. Если \(\Delta_m>0\), пусть \(y_m\) — ближайшая точка
\(\mathcal L_m^X(x,r)\), а \(n_m\) — её первый отделяющий уровень в
блоке. Обозначим через \(\widehat{\mathfrak c}_m\) минимальную цену
контактного разреза, отделяющего корневую компоненту уровня \(n_m\) от
компоненты, содержащей \(y_m\). Последняя является \(X\)-содержащей,
поэтому любой разрез, допустимый в определении
\(\operatorname{cap}_{n_m,m}\), допустим и для этой задачи. Следовательно,
\begin{equation}
 \widehat{\mathfrak c}_m
 \le\operatorname{cap}_{n_m,m}
 \le\mathfrak c_m.
 \label{eq:target-cut-block-bound}
\end{equation}

Для фильтрации Мандельброта берём \(Y_n=O_n\),
\(P_n(c)=p_n(c)/2\), \(X=\M\). При выборе центров ячеек мера \(\Xi_m\)
есть контактная мера блоковой сетки. Для блоков с
\(\Delta_m>0\) и достаточно больших \(m\) имеем
\begin{equation}
 a_{c,r}\Delta_m\le
 \operatorname{Mod}_{1,m}(\Gamma_m)
 \le\widehat{\mathfrak c}_m
 \le\mathfrak c_m
 \le\Xi_m(\C)
 \le B_{c,r} R_m\epsilon_m,
\label{eq:mandelbrot-hypothesis-chain}
\end{equation}
где \(\epsilon_m=\epsilon_m(r)\). Первое неравенство — \(\mathsf{AM}\);
второе следует из того, что каждый допустимый графовый разрез
пересекает все пути из \(\Gamma_m\), а третье — из
\eqref{eq:target-cut-block-bound}. Четвёртое есть
\eqref{eq:block-capacity-contact-mass}, последнее — \(\mathsf{GC}\).
Поэтому
\[
 \Delta_m\le a_{c,r}^{-1}\mathfrak c_m,\qquad
 \mathfrak c_m\le B_{c,r}R_m\epsilon_m.
\]
При \(\Delta_m=0\) первая оценка тривиальна.
В теореме~\ref{thm:abstract-summable-cut} берём
\(A=a_{c,r}^{-1}\) и \(B=B_{c,r}\). Это именно её две количественные
гипотезы; гипотезы
\(\mathsf{AM}\) и \(\mathsf{GC}\), если они выполнены одновременно,
завершают условный переход к факторизации.
Суммируемость \(\epsilon_m\) доказана в
предложении~\ref{prop:mandelbrot-green-block-scales}; открытыми
остаются именно нижняя оценка модуля и Carleson-оценка полной
контактной массы.

\begin{definition}[Совместная гипотеза для факторизации]
\label{def:factorization-open-hypotheses}
Зафиксируем последовательности \((s_m)\), \((\lambda_m)\) из
предложения~\ref{prop:mandelbrot-green-block-scales}. Для фиксированных
\(c\in\partial\M\), \(r>0\) обозначим через \(\mathsf{AM}(c,r)\) и
\(\mathsf{GC}(c,r)\) соответственно утверждения
\begin{equation}
 \begin{aligned}
 \mathsf{AM}(c,r):\quad&
 \exists a_{c,r}>0\ \exists m_A\in\{2,3,\ldots\}\
 \forall m\ge m_A,\quad
 \Delta_m>0\Longrightarrow
 \operatorname{Mod}_{1,m}(\Gamma_m)\ge a_{c,r}\Delta_m,\\
 \mathsf{GC}(c,r):\quad&
 \exists B_{c,r}>0\ \exists m_G\in\{2,3,\ldots\}\
 \forall m\ge m_G,\quad
 R_m>0\ \land\
 \Xi_m(\C)\le B_{c,r}R_m\epsilon_m.
 \end{aligned}
\end{equation}
Глобальная гипотеза \(\mathsf H_{\M}\) означает, что для каждого
\(c\in\partial\M\) и \(\varepsilon>0\) существует \(0<r<\varepsilon\),
для которого одновременно выполнены \(\mathsf{AM}(c,r)\) и
\(\mathsf{GC}(c,r)\). Для фиксированных блоковых масштабов
\(\sum_m\epsilon_m(r)<\infty\).
\end{definition}

\section{Условная теорема о факторизации локальных компонент}
\label{sec:remaining-factorization}

В этом разделе \(\N=\{0,1,2,\ldots\}\), а
\(\overline B(a,\rho):=\{z\in\C:|z-a|\le\rho\}\) обозначает замкнутый
диск. Положим
\[
 K:=\overline B(0,2),\qquad f_c(z):=z^2+c,\qquad
 p_0(c):=0,\qquad p_{j+1}(c):=p_j(c)^2+c,
\]
\[
 \M:=\{c\in\C:\sup_{j\in\N}|p_j(c)|<\infty\},\qquad
 O_N:=\{c\in K:p_j(c)\in K\text{ для }0\le j\le N\}.
\]
Это семейство множеств соответствует радиусу выхода \(2\):
если \(|c|\le2\) и \(|z|>2\), то
\(|f_c(z)|\ge|z|^2-2>|z|\); если \(|c|>2\), то начиная с
\(|p_1(c)|=|c|\) имеем
\(|p_{j+1}(c)|\ge|p_j(c)|^2-|c|\ge|p_j(c)|^2-|p_j(c)|>|p_j(c)|\).
Поэтому \(\M=\bigcap_{N\in\N}O_N\). Из определения также следует
вложенность \(O_L\subseteq O_N\) при \(L\ge N\).
Подмножества \(\C\) рассматриваются с евклидовой подпространственной
топологией. Для \(Y\subseteq\C\) через \(\pi_0(Y)\) обозначим множество
его связных компонент, то есть максимальных связных подмножеств \(Y\).
При \(c\in Y\) через \(\operatorname{Comp}_c(Y)\) обозначим компоненту,
содержащую \(c\).
Для \(c\in\M\), \(r>0\) и \(N\in\N\) положим
\[
 F_N(c,r):=O_N\cap\overline B(c,r),\qquad
 C_N(c,r):=\operatorname{Comp}_c(F_N(c,r)).
\]
Иными словами, \(C_N(c,r)\) — связная компонента внешнего среза
\(F_N(c,r)\), содержащая выделенную точку \(c\).
Так как \(c\in\M=\bigcap_NO_N\), имеем \(c\in F_N(c,r)\), поэтому
\(C_N(c,r)\) определена.
Если \(L\ge N\) и \(0<\delta<r\), то
\[
 X_L(c,\delta):=O_L\cap\overline B(c,\delta)
 \subseteq F_N(c,r).
\]
\noindent
Будем называть \(F_N(c,r)\) локальным срезом внешнего приближения
\(O_N\), а \(X_L(c,\delta)\) — вложенным срезом глубины \(L\).
Обозначим через \(\iota_{L,N}^{c;\delta,r}:X_L(c,\delta)\hookrightarrow
F_N(c,r)\) это включение. Оно задаёт отображение компонент
\[
 \varphi_{L,N}^{c;\delta,r}:=
 \pi_0\!\left(\iota_{L,N}^{c;\delta,r}\right):
 \pi_0\bigl(X_L(c,\delta)\bigr)\longrightarrow
 \pi_0\bigl(F_N(c,r)\bigr),
\]
которое каждой компоненте вложенного среза ставит в соответствие
единственную компоненту среза \(F_N(c,r)\), содержащую её.

\begin{theorem}[Условная теорема о факторизации локальных компонент]
\label{thm:remaining-component-factorization}
Предположим, что для каждого \(c\in\partial\M\) и
\(\varepsilon>0\) существуют общие \(0<r<\varepsilon\) и допустимые
блочные данные, для которых одновременно выполнены открытые оценки
\(\mathsf{AM}(c,r)\) и \(\mathsf{GC}(c,r)\) (совместная гипотеза
\(\mathsf H_{\M}\)). Тогда для введённой системы внешних множеств
\((O_N)\) выполнено условие \(\mathsf{UC}\), то есть
\begin{equation}
 \begin{gathered}
 \forall c\in\M\ \forall\varepsilon>0\
 \exists r,\delta\in\R,\quad 0<\delta<r<\varepsilon,\qquad\\
 \forall N\in\N\ \exists L\in\N,\quad N\le L,\qquad
 \forall A\in\pi_0\bigl(O_L\cap\overline B(c,\delta)\bigr),\\
 \varphi_{L,N}^{c;\delta,r}(A)=C_N(c,r).
 \end{gathered}
 \label{eq:factorization-target}
\end{equation}
\end{theorem}
\begin{proof}
Зафиксируем \(c\in\M\) и \(\varepsilon>0\). Если
\(c\in\operatorname{int}\M\), выберем \(0<r<\varepsilon\) так, чтобы
\(\overline B(c,r)\subseteq\M\), и положим \(\delta=r/2\). Тогда
\(F_N(c,r)=\overline B(c,r)\) для всех \(N\), а
\(O_N\cap\overline B(c,\delta)=\overline B(c,\delta)\); можно взять
\(L=N\).

Пусть теперь \(c\in\partial\M\).
\begin{fuzzystep}
\noindent\textbf{Используемые открытые оценки:}
\(\mathsf{AM}(c,r)\) и \(\mathsf{GC}(c,r)\).
Глобальное предположение теоремы обеспечивает общие \(r\), блочные
данные и константы \(a_{c,r},B_{c,r}>0\), для которых они выполнены.
Цепочка~\eqref{eq:mandelbrot-hypothesis-chain} даёт для всех
достаточно больших \(m\)
\[
 \Delta_m\le a_{c,r}^{-1}\mathfrak c_m,\qquad
 \mathfrak c_m\le B_{c,r}R_m(c,r)\epsilon_m.
\]
При \(\Delta_m=0\) первое неравенство тривиально. Кроме того,
\(\sum_m\epsilon_m<\infty\) по
предложению~\ref{prop:mandelbrot-green-block-scales}, \(R_m>0\)
начиная с некоторого индекса по \(\mathsf{GC}\), а множества
\(O_N\cap\overline B(c,r)\) имеют конечное число компонент по
полуалгебраичности. Поэтому
теорема~\ref{thm:abstract-summable-cut} применима с
\(A=a_{c,r}^{-1}\), \(B=B_{c,r}\), \(Z=K\), \(X=\M\) и \(Y_N=O_N\).
После отбрасывания конечного числа блоков она даёт \(0<\delta<r\) и
для каждого \(N\) индекс \(L\ge N\), для которого
\[
 O_L\cap\overline B(c,\delta)\subseteq C_N(c,r).
\]
\end{fuzzystep}
Следовательно, каждая компонента этого вложенного среза отображается
в \(C_N(c,r)\), что и есть \eqref{eq:factorization-target}.
\end{proof}

\noindent\emph{Смысл формулировки.}
Индекс \(N\) задаёт последний проверяемый член орбиты. При переходе
от \(N\) к \(N+1\) добавляется условие \(p_{N+1}(c)\in K\), поэтому
\(O_{N+1}=O_N\cap p_{N+1}^{-1}(K)\subseteq O_N\). С ростом \(N\)
множества \(O_N\) уточняют конечные внешние приближения \(\M\).

Для каждого \(c\in\M\) и \(\varepsilon>0\) выбираются общие для всех
\(N\) радиусы \(0<\delta<r<\varepsilon\), а для каждого \(N\) —
индекс \(L\ge N\). Обозначим
\[
 X_L:=X_L(c,\delta),\qquad F_N:=F_N(c,r).
\]
Вложенный срез включён в срез \(F_N\), поэтому каждая его компонента
\(A\in\pi_0(X_L)\) содержится в единственной компоненте \(F_N\).
Теорема утверждает, что при подходящем \(L\) эта компонента для каждого
\(A\) равна \(C_N(c,r)\), то есть \(X_L\subseteq C_N(c,r)\).

Срезы на рисунках представлены сеточными выборками.\footnote{Сеточная
выборка \(E\subseteq\C\) —
растр конечной прямоугольной сетки параметров \(c=x+iy\), в котором
пиксель окрашен тогда и только тогда, когда соответствующий узел
принадлежит \(E\).}
Для конкретного примера выбрана граничная точка \(c_0=\frac14\) —
острие главной кардиоиды \(\M\). Сколь угодно близко к \(c_0\) есть
вещественные параметры вне \(\M\), поэтому здесь особенно важно
проследить, как конечные внешние приближения уточняют локальную
геометрию у границы.

\begin{figure}[H]
 \begin{minipage}[t]{0.40\textwidth}
  \vspace{0pt}
  \centering
  \includegraphics[width=\linewidth]{docs/figures/factorization-critical-orbit.pdf}
  \captionof{figure}{Критическая орбита при граничном параметре \(c_0=\frac14\).}
  \label{fig:factorization-orbit}
 \end{minipage}\hfill
 \begin{minipage}[t]{0.56\textwidth}
  \vspace{0pt}
  \scriptsize
  Для \(c_0=\frac14\) имеем \(p_0(c_0)=0\),
  \(p_1(c_0)=\frac14\), \(p_2(c_0)=\frac5{16}\), а далее
  \(p_{n+1}(c_0)=p_n(c_0)^2+\frac14\). Орбита возрастает к
  параболической неподвижной точке \(\frac12\), поэтому все её члены
  лежат в \(K=\overline B(0,2)\) и \(c_0\in\M\).
  Для вещественного \(c>\frac14\) орбита строго возрастает; если бы она
  была ограничена, её предел был бы вещественной неподвижной точкой
  \(x^2+c=x\), которой при \(c>\frac14\) нет. Значит, такие \(c\) не
  принадлежат \(\M\), и \(c_0\in\partial\M\).
 \end{minipage}
\end{figure}

\begin{figure}[H]
 \begin{minipage}[t]{0.40\textwidth}
  \vspace{0pt}
  \centering
  \includegraphics[width=\linewidth]{docs/figures/factorization-global-levels.pdf}
  \captionof{figure}{Внешние приближения и окна у точки \(c_0=\frac14\in\partial\M\).}
  \label{fig:factorization-global}
 \end{minipage}\hfill
 \begin{minipage}[t]{0.56\textwidth}
  \vspace{0pt}
  \scriptsize
  На рис.~\ref{fig:factorization-global} взяты
  \(0<\delta=\frac15<r=\frac25<\varepsilon=\frac12\), \(N=2\), \(L=20\).
  Условие \(O_j\) проверяет орбиту до индекса \(j\), поэтому
  \(\M\subseteq O_{20}\subseteq O_2\).
  \[
   \begin{aligned}
    F_2&=O_2\cap\overline B(c_0,r),\\
    X_{20}&=O_{20}\cap\overline B(c_0,\delta)\subseteq F_2.
   \end{aligned}
  \]
  На фоне отдельно показаны слои \(O_2\setminus O_{20}\) и \(O_{20}\);
  пунктиром выделена граница \(O_{20}\). Оранжевая заливка показывает не
  всё \(O_{20}\), а только локальный срез \(X_{20}\).
  Полупрозрачная голубая заливка выделяет \(F_2\), оранжевая — \(X_{20}\);
  подписи нанесены прямо на них. Красная штриховая и
  оранжевая точечная окружности ограничивают круги радиусов \(r\) и
  \(\delta\). Включение \(X_{20}\hookrightarrow F_2\) задаёт отображение
  \(\varphi_{20,2}^{c_0;\delta,r}\): каждая компонента \(X_{20}\)
  относится к компоненте \(F_2\), содержащей её. Здесь это
  \(C_2(c_0,r)\), как
  показано на рис.~\ref{fig:factorization-target}.
  Звезда отмечает граничный параметр \(c_0=\frac14\), центр обоих кругов.
 \end{minipage}
\end{figure}

\begin{figure}[H]
 \begin{minipage}[t]{0.40\textwidth}
  \vspace{0pt}
  \centering
  \includegraphics[width=\linewidth]{docs/figures/factorization-target.pdf}
  \captionof{figure}{Срезы \(F_2\), \(X_{20}\) и граница \(O_{20}\) у граничной точки \(c_0=\frac14\).}
  \label{fig:factorization-target}
 \end{minipage}\hfill
 \begin{minipage}[t]{0.56\textwidth}
  \vspace{0pt}
  \scriptsize
  Для \(c=c_0+u=\frac14+u\), \(|u|\le\frac25\), имеем оценки
  \[
   |c|\le\frac{13}{20}<2,\qquad
   |p_2(c)|=|c(c+1)|
   \le\frac{13}{20}\cdot\frac{33}{20}
   =\frac{429}{400}<2.
  \]
  Поэтому \(F_2=\overline B(c_0,\frac25)=C_2(c_0,\frac25)\).
  На рис.~\ref{fig:factorization-target} синяя заливка показывает
  сеточную выборку \(F_2\), оранжевая заливка и контур — срез \(X_{20}\),
  а зелёный штрихпунктир — границу \(O_{20}\). Метка \(A\) указывает на
  выбранную компоненту \(A\in\pi_0(X_{20})\), звезда — граничный
  параметр \(c_0\).
  В этом примере включение задаёт отображение компонент
  \[
   \begin{aligned}
    X_{20}&\subseteq F_2=C_2(c_0,r),\\
    \varphi_{20,2}^{c_0;\delta,r}&:
       \pi_0(X_{20})\to\pi_0(F_2),\\
    A&\longmapsto C_2(c_0,r).
   \end{aligned}
  \]
  Это один численный случай условия \(\mathsf{UC}\): при \(N=2\)
  выбран \(L=20\), а образом каждой компоненты вложенного среза служит
  \(C_2(c_0,r)\).
 \end{minipage}
\end{figure}

\noindent\emph{Геометрический смысл теоремы.}
Рис.~\ref{fig:factorization-global} и~\ref{fig:factorization-target}
показывают срезы \(X_{20}\subseteq F_2\) у граничной точки
\(c_0=\frac14\) и выбранную компоненту \(A\in\pi_0(X_{20})\), которая
содержится в \(C_2(c_0,r)\). Теорема
распространяет это свойство на каждую компоненту \(X_L\) при достаточной
глубине \(L\):
\[
 X_L(c,\delta)\subseteq C_N(c,r),\qquad
 \varphi_{L,N}^{c;\delta,r}(A)=C_N(c,r)
 \quad\text{для каждого }A\in\pi_0(X_L(c,\delta)).
\]
У среза \(X_L\) могут быть разные компоненты, и все они переходят в
одну компоненту \(C_N(c,r)\). Поэтому
\[
 \operatorname{im}\varphi_{L,N}^{c;\delta,r}
 =\{C_N(c,r)\}.
\]
На множестве компонент \(\pi_0(X_L(c,\delta))\) зададим отношение
эквивалентности
\[
A\sim B
\quad\Longleftrightarrow\quad
\varphi_{L,N}^{c;\delta,r}(A)=\varphi_{L,N}^{c;\delta,r}(B).
\]
Фактор-множество \(\pi_0(X_L(c,\delta))/\sim\) состоит из классов
компонент с одинаковым образом. Множество \(X_L(c,\delta)\) содержит
\(c\), поэтому у него есть компонента \(A\). По теореме все компоненты
эквивалентны, и фактор-множество состоит из одного класса:
\[
\pi_0(X_L(c,\delta))/\sim=\{[A]\}.
\]
Элементами фактор-множества служат классы связных компонент \(X_L\).
На правой части рис.~\ref{fig:factorization-component-geometry}
области \(A_0'\) и \(A_1'\) остаются раздельными; соотношение
\(A_0'\sim A_1'\) означает, что их образы совпадают и обе компоненты
входят в один класс \([A]\).
\[
\begin{aligned}
 q &: \pi_0(X_L(c,\delta))\longrightarrow
      \pi_0(X_L(c,\delta))/\sim,
      &q(A)&=[A],\\
 \bar\varphi_{L,N}^{c;\delta,r}&:
      \pi_0(X_L(c,\delta))/\sim\longrightarrow\pi_0(F_N(c,r)),
      &\bar\varphi_{L,N}^{c;\delta,r}([A])&=C_N(c,r),\\
 \varphi_{L,N}^{c;\delta,r}&=\bar\varphi_{L,N}^{c;\delta,r}\circ q.
\end{aligned}
\]
Таким образом, одноэлементный объект здесь — фактор-множество по
\(\sim\); его единственный класс соответствует единственному значению
\(\{C_N(c,r)\}\) отображения компонент.
\(\mathcal R_N:=\pi_0(F_N(c,r))\setminus\{C_N(c,r)\}\) обозначает
остальные компоненты \(F_N\).
Для изображения ниже обозначим через \(D\in\mathcal R_N\) одну такую
компоненту.
Для наглядности рис.~\ref{fig:factorization-component-geometry}
сравнивает промежуточную глубину \(\ell\) и достаточную глубину \(L\),
где \(N\le\ell<L\). В изображённом случае срез \(X_\ell\) имеет
компоненты \(A_0,A_1\subseteq C_N(c,r)\) и \(B\subseteq D\), где
\(D\in\mathcal R_N\); поэтому
\(\operatorname{im}\varphi_{\ell,N}^{c;\delta,r}=\{C_N(c,r),D\}\).
Глубина \(L\) задаёт срез \(X_L\subseteq X_\ell\), в котором остаются
компоненты \(A_0',A_1'\) внутри \(C_N(c,r)\). Поскольку \(D\) отделена
от \(C_N(c,r)\), имеем \(B\cap X_L=\varnothing\), и
\(\operatorname{im}\varphi_{L,N}^{c;\delta,r}=\{C_N(c,r)\}\).
Так рисунок показывает, как дополнительное условие глубины сужает
множество значений отображения. Число и форма изображённых компонент
условны. Точка \(c\) на схеме расположена у края \(C_N(c,r)\), чтобы
связать её с граничным численным примером выше; это расположение также
схематично.

\begin{figure}[H]
 \centering
 \includegraphics[width=0.96\textwidth]{docs/figures/factorization-component-geometry.pdf}
 \captionof{figure}{Срезы \(X_\ell\supseteq X_L\): раздельные компоненты \(A_i'\) образуют класс \([A]\), соответствующий \(C_N(c,r)\).}
 \label{fig:factorization-component-geometry}
\end{figure}
